\begin{proposition}[DP37: the true dual torus and a primitive normalized kernel]\label{DP37}
In the actual model of DP32 over $\mathcal O=W(\F_q)$, the automorphism
\[
 A_\zeta(\mathbf X)=(\zeta^pX_1,\ldots,\zeta^{p^{h-1}}X_{h-1}),\qquad\zeta^{q-1}=1,
\]
satisfies $A_\zeta^D=[\zeta]_{G_h}$. Consequently the primitive formal Poincar\'e biextension has an integral normalized frame invariant under
\[
 ([\zeta],A_\zeta^{-1}).
\]
In that frame there is a unique invariant rational logarithmic splitting $\beta$. It vanishes on both axes, has no bidegree-$(1,1)$ term, and has the integral exponential coboundaries required in DP36. Thus for odd $p$ DP36 applies to an actual primitive formal Poincar\'e representative. Its compatible splitting still has the denominators proved in DP35. This does not identify the particular section, gauge, path or harmonic operators of a global theta correspondence.
\end{proposition}
\begin{proof}
The explicit monomial map $M$ of DP32 commutes with $A_\zeta$: the ordinary arrows use $p\cdot p^{i-1}=p^i$, and the exceptional arrow uses $p^{h+1}\equiv p\pmod{q-1}$. Hence its logarithm and integral group law commute as well. This proves that $A_\zeta$ is an integral automorphism. It suffices first to take $\zeta$ primitive. Put $a=A_\zeta^D$ and write $\lambda=da$ on the one-dimensional tangent space of $G_h$.

Tate's functorial Hodge--Tate sequence \cite[\S4, Corollary 2 following Theorem 3, p.~180]{Tate67}, over the completed algebraic closure $C$ of $\operatorname{Frac}(\mathcal O)$, is
\[
 0\longrightarrow \Lie(H)_C
 \longrightarrow \Hom_{\Z_p}(T_pG_h,C)
 \longrightarrow \Lie(G_h)_C^*\otimes_C\Hom_{\Z_p}(\Z_p(1),C)
 \longrightarrow0.
\]
On its middle term $a$ acts by precomposition, on the first term by $d(a^D)=dA_\zeta$, and on the cotangent factor by pullback, with scalar $\lambda$. These are adjoint actions, so neither Hom nor the cotangent introduces an inverse eigenvalue. This also follows directly by differentiating the natural Cartier identity $\langle aP,Q\rangle=\langle P,a^DQ\rangle$. The cyclotomic factor is unchanged.

The characteristic polynomial $P_a(Z)$ of $a$ on the rank-$h$ $\Z_p$ Tate module is consequently
\[
 P_a(Z)=(Z-\lambda)\prod_{i=1}^{h-1}(Z-\zeta^{p^i}).
\]
It is monic of degree $h$ over $\Q_p$. The primitive root $\zeta^p$ has degree exactly $h$: its Frobenius conjugates are $\zeta^{p^j}$, and $p^m\not\equiv1\pmod{p^h-1}$ for $0<m<h$. Therefore its degree-$h$ minimal polynomial divides and equals $P_a$. Comparing the one missing root proves $\lambda=\zeta$. In characteristic zero a one-dimensional formal endomorphism is determined by its differential: conjugate by the normalized logarithm and compare the coefficient of $X^{n-1}Y$ in additivity to kill all degrees $n\ge2$. Thus $a=[\zeta]$. Taking powers proves the equality for every root, and its integral identity persists under base change.

The evaluation pairing giving the primitive Poincar\'e biextension is preserved by $([\zeta],A_\zeta^{-1})$. Functoriality of the pairing construction, not just equality of isomorphism classes, gives a compatible action of the finite group $\Gamma=\mu_{q-1}$ on the biextension and its axial rigidifications. Choose an integral normalized frame $e$ as in DP35 and write $g(e)=c_g e$. The units $c_g$ are one on both axes and satisfy $c_{\sigma g}=\sigma(c_g)c_\sigma$. Set $N=q-1$ and
\[
 u=\left(\prod_{g\in\Gamma}c_g\right)^{1/N}.
\]
The unique root with constant term one is integral: coefficientwise recursion divides only by the unit $N$. It is one on both axes. The cocycle identity gives
\[
 \sigma(u)=u/c_\sigma,
\]
so $e'=ue$ is an invariant integral normalized frame. Its two addition cocycles are invariant integral units.

Use DP35 to obtain a rational logarithmic splitting $\beta_0$ of these cocycles in the additive logarithm coordinates, and average
\[
 \beta=\frac1N\sum_{g\in\Gamma}g\beta_0.
\]
Both additive coboundaries remain the logarithms of the same invariant integral cocycles. The axes remain zero. The characters are exactly $(\zeta,\zeta^{-p},\ldots,\zeta^{-p^{h-1}})$ by the first part of the proof. No bilinear monomial $xy_i$ is invariant, because $1-p^i$ is not divisible by $q-1$ for $1\le i<h$. Thus $\beta$ has no bidegree-$(1,1)$ term. Any two rational splittings differ by a bilinear form by DP35, so the same character observation proves uniqueness among invariant splittings. All hypotheses of DP36 have now been obtained for the primitive class. Its nontriviality and denominator conclusion are unchanged.
\end{proof}
